\section{Proofs for the high-dimensional mean-field approximation}

\label{2302:sec:app:hd_mean_field}
\subsection{Preliminaries}
We first provide several bounds on expectations of functions of Gaussians, that will be used throughout this section. We begin by recalling the classical Gaussian Poincaré inequality (see e.g. \cite{boucheron2013concentration}, Theorem 3.20):
\begin{lemma}\label{2302:lem:gauss_poincare}
    Let $f: \mathbb{R}^d \to \mathbb{R}$ be a differentiable function, and $X \sim \mathcal N(0, I)$. Then
    \[ \mathrm{Var}(f(X)) \leq \mathbb E\left[\lVert\nabla f(X) \rVert^2 \right]. \]
    If $X \sim \mathcal{N}(0, \Sigma)$, the following bounds hold instead:
    \[ \mathrm{Var}(f(X)) \leq \mathbb E\left[\langle\nabla f(X), \Sigma \nabla f(X) \rangle \right] \leq \lVert \Sigma \rVert_{\mathrm{op}} \cdot \mathbb E\left[\lVert\nabla f(X) \rVert^2 \right] . \]
\end{lemma}

Now, we provide some concentration bounds for expectations of random variables. Let $h:\mathbb{R}\to\mathbb{R}$ be a function, and $z$ a random variable; our goal is to bound
\[ \left|\mathbb{E}[h(z)] - h(\mathbb{E}[z])\right| \]
under different assumptions on $f, z$.
\begin{lemma}\label{2302:lem:average_lipschitz_function}
    Assume that $h$ is $L$-Lipschitz, and that $z$ has a second moment. Then
\begin{equation}
        \left|\mathbb{E}[h(z)] - h(\mathbb{E}[z])\right| \leq L\sqrt{\Var(z)}
    \end{equation}
\end{lemma}
\begin{proof}
    By Jensen's inequality, we have
    \begin{align*}
        \left|\mathbb{E}[h(z)] - h(\mathbb{E}[z])\right| &\leq \mathbb{E}\left[|h(z) - h(\mathbb{E}[z])| \right] \\
        &\leq L \mathbb{E}[|z - \mathbb{E}[z]|] \\
        &\leq L\sqrt{\Var(z)}
    \end{align*}
    where the last line uses the Cauchy-Schwarz inequality.
\end{proof}

\begin{lemma}\label{2302:lem:average_twice_diff_function}
    Assume that $h$ is twice differentiable, with $\norm{h^{(k)}}_\infty \leq K$, and that $z$ has a fourth moment. Then
    \begin{equation}
        \left|\mathbb{E}[h(z)] - h(\mathbb{E}[z])\right| \leq \frac K2 \sqrt{\mathbb{E}\left[(z-s)^4 \right]} 
    \end{equation}
\end{lemma}
\begin{proof}
    For simplicity, let $s = \mathbb{E}[z]$. By the Lagrange formula for Taylor series, we can write
    \begin{equation}
        h(x) = h(s) + (x-s)h'(s) + (x-s)^2 R(x)
    \end{equation}
    where $R$ is bounded by $K/2$. Plugging $z$ into this equation,
    \begin{equation}
        \left|\mathbb{E}[h(z)] - h(s)\right| = \left|\mathbb{E}\left[ (z-s)^2 R(z) \right] \right|
        \leq \sqrt{\mathbb{E}\left[(z-s)^4 \right] \mathbb{E}\left[ R(z)^2 \right]}
    \end{equation}
    via the Cauchy-Schwarz inequality, and the result ensues.
\end{proof}

\subsection{Proof of Lemma~\ref{2302:lem:Psi_lipschitz}}

We only show the result for $\Psi^{(M)}$; the one for $\Psi^\bot$ ensues from similar methods. Recalling the expansion  in \eqref{2302:eq:app:ss_ode_expansion} as a sum of 3-point correlation functions, we define the following function of $3 \times 3$ matrices:
\begin{equation}
    f(\Sigma) = \mathbb{E}_{\vec{z} \sim \mathcal N(\vec{0}, \Sigma)}\left[ \sigma'(z_1)z_2\sigma(z_3). \right]
\end{equation}
Then $\Psi^{(M)}(\Omega)$ is simply an average of $f(\Sigma)$ for $\Sigma$ submatrices of $\Omega$. We first show the following bound:
\begin{lemma}\label{2302:lem:psi_reduced_lipschitz}
    Under Assumption~\ref{2302:assump:activation}, the function $f$ satisfies the following inequality: for any $\Sigma$, $\Sigma'$ such that $\norm{\Sigma' - \Sigma}_\infty \leq K$,
    \begin{equation}
        \left|f(\Sigma') - f(\Sigma)\right| \leq C(1 + \sqrt{\Sigma_{22}})\norm{\Sigma' - \Sigma}_\infty
    \end{equation}
\end{lemma}
\begin{proof}
    Let $\Sigma$ and $\Sigma'$ be two positive semidefinite matrices, and $\delta\Sigma = \Sigma' - \Sigma$. We can write $\delta\Sigma = \delta\Sigma_1 - \delta\Sigma_2$ where both $\delta\Sigma_i$ are positive, and
    \[ \norm{\delta\Sigma_i}_{\infty} \leq C\norm{\delta\Sigma_i}_{\mathrm{op}} \leq  C\norm{\delta\Sigma}_{\mathrm{op}} \leq C'\norm{\delta\Sigma}_{\infty} \]
    using norm equivalence. Hence, we shall assume from now on that $\delta\Sigma$ is positive semidefinite, and write
    \[z' = z + \delta z\]
    where $\delta z \sim \mathcal N(\vec{0}, \delta\Sigma)$ is independent from $z$. Define
    \[ \delta\sigma = \sigma(z_3') - \sigma(z_3) \quad \text{and} \quad \delta\sigma' = \sigma'(z_1') - \sigma'(z_1) \]
    Then
    \begin{align*}
        f(\Sigma') - f(\Sigma) = \quad&\mathbb{E}\left[ \sigma'(z_1)z_2\delta\sigma + \sigma'(z_1)\delta z_2 \sigma(z_3) + \delta \sigma' z_2\sigma(z_3) \right] \\
        + \ &\mathbb{E}\left[ \delta\sigma'\delta z_2 \sigma(z_3) + \delta\sigma' z_2 \delta\sigma + \sigma'(z_1)\delta z_2 \delta \sigma \right] \\
        + \ & \mathbb{E}[\delta\sigma' \delta z_2 \delta \sigma]
    \end{align*}
    Since $\delta z$ is independent from $z$, and $\sigma$ satisfies Assumption~\ref{2302:assump:activation}, we can apply Lemma~\ref{2302:lem:average_twice_diff_function} to get
    \[ |\mathbb{E}_{\delta z_3}[\delta\sigma]| \leq C\delta\Sigma_{33}, \]
    and hence by the Cauchy-Schwarz inequality
    \[ \left|\mathbb{E}\left[ \sigma'(z_1)z_2\delta\sigma \right]\right| \leq C' \sqrt{\Sigma_{22}}  \delta\Sigma_{33}.  \]
    A similar bound holds for the two other terms of the first line. For the second line, another application of Cauchy-Schwarz yields
    \begin{align*}
        \mathbb{E}_{\delta z}\left[\delta\sigma' z_2 \delta\sigma\right] &\leq z_2 \sqrt{\mathbb{E}[(\delta\sigma)^2]}\sqrt{\mathbb{E}[(\delta\sigma')^2]} \\
        &\leq z_2 \sqrt{\delta\Sigma_{11}\delta\Sigma_{33}}
    \end{align*}
    by a combination of Lemmas~\ref{2302:lem:gauss_poincare} and \ref{2302:lem:average_lipschitz_function}. A similar bound holds for the third line, and it is easily checked that all of the obtained bounds are lower than the one of Lemma~\ref{2302:lem:psi_reduced_lipschitz}.
\end{proof}

Now, the expansion of $\Psi^{(M)}(\tilde\Omega) - \Psi^{(M)}(\bar\Omega)$ is an average of terms of the form $f(\tilde\Sigma) - f(\bar\Sigma)$, where $(\tilde\Sigma - \bar\Sigma)_{ij}$ is either $0$ or $\sqrt{q_{ii}^\bot q_{jj}^\bot} \xi_{ij}$. Hence, we can write each of those terms as $h(\xi_{ij}) - h(0)$, and by Assumption~\ref{2302:assump:boundedness} and Lemma~\ref{2302:lem:psi_reduced_lipschitz}, the function $h$ is Lipschitz. Lemma~\ref{2302:lem:Psi_lipschitz} then ensues from an application of Lemma~\ref{2302:lem:average_lipschitz_function}.

\subsection{Proof of Equation~\texorpdfstring{\eqref{2302:eq:mf:concentration}}{\ref{2302:eq:mf:concentration}}}
We begin by showing the following lemma. Recall the definition of $\mathcal{E}$:
\begin{equation}\label{2302:eq:app:def_error}
    \mathcal E = \frac1k\sum_{r=1}^k \sigma(\lambda_r^\star) - \frac1p \sum_{i=1}^p \sigma(\lambda_i)
\end{equation}

\begin{lemma}\label{2302:lem:app:gauss_poincare}
    There exists a constant $c \geq 0$ such that for any choice of $\vec{\lambda}^\star$,
    \[ \mathrm{Var}_{\vec{\lambda^\bot}}\left( \mathcal E \right) \leq \frac{c  \lVert Q^\bot \rVert_{\mathrm{op}}}{p} \]
\end{lemma}

\begin{proof}
    We apply the Gauss-Poincaré inequality of Lemma~\ref{2302:lem:gauss_poincare} to
    \[ f(\vec{\lambda}^\bot) = \sum_{i=1}^p \sigma([M P^{-1}\vec{\lambda}^\star]_i + \lambda_i^\bot), \quad \text{which implies} \quad [\nabla f]_i = \sigma'(\lambda_i)\]
    Whenever $\sigma$ is Lipschitz, we thus have $\lVert\nabla f(\vec{\lambda}^\bot) \rVert^2 \leq c p$, and the lemma ensues.
\end{proof}

We are now in a position to show Equation~\eqref{2302:eq:mf:concentration}. For brevity, we denote by $\mathbb E$ (resp. $\mathbb E_{\text{MF}}$) the expectations with respect to $\vec\lambda^\bot \sim \mathcal N(\vec{0}, Q^\bot)$ (resp. $\vec\lambda^\bot \sim \mathcal N(\vec{0}, \mathrm{diag}(Q^\bot))$). Since the marginals of both distributions are the same, and by linearity,
\[ \mathbb E\left[f(\lambda^\bot_i)\right] =  \mathbb E_{\text{MF}}\left[f(\lambda^\bot_i)\right] \quad \text{and} \quad  \mathbb E\left[\mathcal E\right] =  \mathbb E_{\text{MF}}\left[\mathcal E\right]. \]
Under the distribution $\mathcal N(\vec{0}, \mathrm{diag}(Q^\bot))$, $\lambda_i^\bot$ is almost independent from $\mathcal E$, except for the term containing $\sigma(\lambda_i)$. We can thus write, for any $\vec{\lambda}^\star$
\[ \mathbb E_{\text{MF}}\left[\mathcal E_i \lambda_i^\bot \right] = \mathbb E \left[ \sigma'(\lambda_i)\lambda_i^\bot \right] \mathbb E\left[ \mathcal E \right] - \frac1p\mathbb E_{\text{MF}}\left[ \sigma'(\lambda_i)\lambda_i^\bot \sigma(\lambda_i) \right] + \frac1p \mathbb E\left[\sigma'(\lambda_i)\lambda_i^\bot \right] \mathbb E\left[\sigma(\lambda_i) \right]  \]
Hence,
\begin{equation}\label{2302:eq:app:expectation_difference}
\mathbb E\left[\mathcal E_i \lambda_i^\bot \right] - \mathbb E_{\text{MF}}\left[\mathcal E_i \lambda_i^\bot \right] = \mathbb E \left[ \sigma'(\lambda_i)\lambda_i^\bot \left(\mathcal E - \mathbb E\left[ \mathcal E \right]\right)\right]  - \frac1p\mathbb E\left[ \sigma'(\lambda_i)\lambda_i^\bot \left(\sigma(\lambda_i) - \mathbb E[\sigma(\lambda_i)] \right)\right] 
\end{equation} 
Now, using the Cauchy-Schwarz inequality,
\[  \left|\mathbb E \left[ \sigma'(\lambda_i)\lambda_i^\bot \left(\mathcal E - \mathbb E\left[ \mathcal E \right]\right)\right]\right| \leq \sqrt{\mathbb E \left[ \left(\sigma'(\lambda_i)\lambda_i^\bot\right)^2\right]\mathbb E \left[\left(\mathcal E - \mathbb E\left[ \mathcal E \right]\right)^2\right]}.\]
For Lipschitz $\sigma$, the first term is easily bounded by $c  Q_{ii}^\bot$, and the second is exactly the variance computed in Lemma~\ref{2302:lem:app:gauss_poincare}. The second term in \eqref{2302:eq:app:expectation_difference} being clearly negligible before the first, we finally get
\begin{equation}
    \left|\mathbb E\left[\mathcal E_i \lambda_i^\bot \right] - \mathbb E_{\text{MF}}\left[\mathcal E_i \lambda_i^\bot \right]\right| \leq c  \sqrt{Q_{ii}^\bot \cdot \frac{\lVert Q^\bot \rVert_{\mathrm{op}}}{p}},
\end{equation}
and Equation~\eqref{2302:eq:mf:concentration} ensues by taking the expectation w.r.t $\vec{\lambda}^\star$ on both sides.

% %%%%%%%%%%%%%%%%%%%%%%%%%%%%%%%%%%%%%%%%%%%%%%%%%%%%%%%%%%%%%%%%%%%%%%%%%%%%%%%
